\scanpage{47}
\deleted{$a$ à une forme linéaire sur $B$, et $\pi_1$ à la
section $e_B$ de $B$ (au signe près\ldots) cette relation
s'écrit $a(e_B)=-2$, et \deleted{cette équation}
on vérifie (\deleted{\unsure{sachant bien}} \inserted{faisant attention} au \unsure{traquenard|traitement} des
signes) que $a$ correspond à $\underline b=\mathrm{tr}$ \ldots\ Le
seul cas est celui où}

Cette quasi-réflexion est une
réflexion i.e. est $\ne$ identité en chaque fibre, \unsure{sss|ssi}
en tt pt $s\in S$ de car résiduelle 2, on a
$\mathrm{tr}_s$ sur $B_s$ non (id\textsuperscript{t}) nulle. Ceci résulte \unsure{des|de} la

\uline{Proposition} Pour qu'une quasi-réflexion d'un
module \deleted{\unsure{quadratique, soit rg 2}} \inserted{\unsure{quelconque|partout de}} \deleted{de} $B$, (par r. à un
$B'$ tel que $B/B'\simeq\check L$ inv. \uline{donné}) soit l'identité,
il f et s. que $S$ soit de car. 2 (i.e. $2.1_S=0$)
et que la \deleted{\unsure{$f$}} section de $L\otimes B$ qu'elle définit soit
nulle \ldots\ donc dans le cas actuel ($B$ de rg 2,
$B'=\mathcal O_S e_B$, donc $L\otimes B\simeq\check B$, que la forme $\underline b=\mathrm{tr}$
sur $B$ qu'elle définit (et qui la définit, car
$\underline b(e_B)=2$ \unsure{---|\ldots}) soit nulle.

\deleted{Supposons}\inserted{Revenons} maintenant au cas où $\underline b=\mathrm{tr}$ \deleted{\unsure{provient}} \inserted{est la}
\inserted{forme trace} d'une structure d'Algèbre sur $B$ ayant $e_B$ comme
unité. On obtient alors la formule bien connue
\[
u.\sigma u=\unsure{\mathcal N|\mathscr N|N}(u)e_B=q(u)e_B \tag*{(86)}
\]
i.e. \qquad $u(\mathrm{tr}\,u-u)=\mathcal N(u)$ \qquad ou encore
\[
\big|\,u^2-u\,\mathrm{tr}\,u+\mathcal N(u)=0 \tag*{(87)}
\]
% archive sheet number: 46
